\emph{Step 1 (the commutant in the limit).} By the definition of the tower as the chain of unital embeddings $(T_{k}, \iota_{k})$ (\claimref{obj.tower.tower}), every element of its algebraic direct limit, for either parity of $p - q$, lies in some floor and is a combination of finite blades in the ground and plane generators. Conjugation by a generator $g$ acts diagonally on blades, $gB = \pm Bg$, so an element commutes with $g$ exactly when each blade in its support does. For a blade $B$ and a generator $g$ (distinct frame generators anticommute, \claimref{std-algebra.orthogonal-anticommute}): if $g \notin B$ then $gB = (-1)^{|B|}Bg$; if $g \in B$ then $gB = (-1)^{|B|-1}Bg$. A blade involves finitely many planes, so some plane generator lies outside it, and commuting with that one forces $|B|$ even; then any plane generator inside $B$ anticommutes with it. So the blades commuting with all of $B_{\infty}$ are the even blades in the ground generators alone, which span $\Cl^{0}(p,q)$. (For $p - q$ even, Theorem \claimref{tower.algebraic-limit-even-charge} identifies the limit further; that is not needed here.)

\emph{Step 2 (the single floor).} On $T_{k}$ the commutant of $B_{k}$ is $A_{k}$, generated by $x^{(k)} = x\,\Omega_{k}$, and $A_{k} \cong \Cl(p,q)$ (Theorem \claimref{tower.commuting-copies}). The element $\Omega_{k} = \omega_{1} \cdots \omega_{k}$ is a product of bivectors of the first $k$ planes (Lemma \claimref{tower.volume-elements-planes}), each of which commutes with $e_{k+1}$ and $f_{k+1}$, while $x$ anticommutes with them; so $x^{(k)}$ anticommutes with the new plane, and of $A_{k}$ only the span of products of an even number of twisted generators, which is $\Cl^{0}(p,q)$ ($\Omega_{k}^{2} = 1$ and $\Omega_{k}$ commutes with $x$), commutes with $T_{k+1}$'s planes.

\emph{Step 3 (what the commutant records).} For $n$ odd the centre of $\Cl^{0}(p,q)$ is $\R$; for $n$ even it is $\mathrm{span}\{1, \omega\}$ with $\omega$ the ground volume element, even and central in $\Cl^{0}$, and $\omega^{2} = (-1)^{n(n-1)/2}(-1)^{q} = (-1)^{r/2}$. Further, $\Cl^{0}(p,q) \cong \Cl(p, q-1)$ (or $\Cl(q, p-1)$) through $e_{i} \mapsto e_{i}e_{n}$, whose type the classification of Theorem \claimref{transport.classification-real-clifford-algebras} gives as listed in (ii), the types that Theorem \claimref{tower.rigidity-clock-frozen} carries unchanged along the tower; the five types are distinct, so with $n$ they fix the class $\{r, -r\}$.

\emph{Step 4 (the reversal).} Reverse every square on one frame. The elements $u_{i} = e_{i}e_{n}$, $i < n$, satisfy $u_{i}^{2} = -Q(e_{i})Q(e_{n})$ and anticommute pairwise, in both signatures, and $Q(e_{i})Q(e_{n})$ is unchanged when both factors change sign. So $\Cl^{0}(p,q)$ and $\Cl^{0}(q,p)$ have one presentation and are isomorphic. The check verifies Steps 1 and 2 on small floors (the commutant of all planes of $T_{k+1}$, restricted to $T_{k}$, is the span of the even ground blades; the commutant on $T_{k}$ alone has dimension $2^{n}$), the common presentation of Step 4, and the centre and $\omega^{2}$ of Step 3, in exact arithmetic.
